% ============================================================
% Why the Golden Ratio Selects the Prime Three
% ============================================================

\documentclass[12pt]{amsart}

\usepackage{amsmath,amssymb,amsthm}
\usepackage{booktabs}
\usepackage{microtype}
\usepackage{url}

\newtheorem{theorem}{Theorem}
\newtheorem{lemma}[theorem]{Lemma}
\newtheorem{proposition}[theorem]{Proposition}
\newtheorem{corollary}[theorem]{Corollary}
\theoremstyle{definition}
\newtheorem{definition}[theorem]{Definition}
\newtheorem{conjecture}[theorem]{Conjecture}
\theoremstyle{remark}
\newtheorem{remark}[theorem]{Remark}

\title{Why the Golden Ratio Selects the Prime Three}

\author{Alexander S.\ Petty}

\email{alexander.petty@gmail.com}
\date{January 2020 (revised April 2026)}
\subjclass[2020]{11A63, 11B83}

\begin{document}
\begin{abstract}
For a prime $p$ and a $b$-smooth integer $m$, the alignment
proportion $\alpha(pm) = (2m{-}1)/(pm{-}1)$ measures the fraction
of the fractional field $\{k/(pm) : 1 \le k \le pm{-}1\}$ whose
repetend matches that of $1/(pm)$. The formula's algebraic
properties are independent of the choice of base.

We prove that $\alpha \ge 1/\varphi$ (the reciprocal of the golden
ratio) if and only if $m \ge \varphi^{2} = \varphi + 1$, provided
$p = 3$. We then establish a stronger result. The self-referential
condition $m^{*}(\tau) = 1/\tau^{2}$ yields a cubic equation that
admits a root in $(0,1)$ for every prime $p \ge 3$, but that factors
over $\mathbb{Q}$ only for $p \in \{2, 3, 5\}$. The two non-trivial
factorizations produce thresholds $1/\varphi$ and $1/\varphi^{2}$,
both in the golden ratio's number field $\mathbb{Q}(\sqrt{5})$. The
golden ratio's minimal polynomial $\tau^{2} + \tau - 1$ divides the
cubic if and only if $p = 3$, with coefficient~$3$ produced by
polynomial division. For $p \ge 7$, the cubic is irreducible and the
self-referential threshold is a cubic irrational outside any quadratic
field.

The concrete setting is decimal arithmetic, where the alignment formula
describes the coherence of fractional fields $\{k/n : 1 \le k \le
n{-}1\}$. But the central results, the golden threshold, the
self-referential characterization, and the algebraic selection of the
prime~$3$, hold for any positional base $b$ with
$b \equiv 1 \pmod{3}$ (including base~$10$).
\end{abstract}
\maketitle

% ============================================================
\section{Introduction}

The decimal expansion of a rational number $k/n$ either terminates or
eventually repeats. When it repeats, the shortest repeating block is
called the \emph{repetend}. The length of the repetend of $1/n$, for
$\gcd(n,10)=1$, equals the multiplicative order
$\operatorname{ord}_{n}(10)$, and the repetends of the family
$\{k/n : \gcd(k,n)=1\}$ are cyclic permutations of one another. These
facts are classical~\cite{hardy}.

Less studied is the following question. Given a fixed denominator $n$,
consider the \emph{full fractional field}
$\mathcal{F}(n) = \{k/n : 1 \le k \le n-1\}$ and ask: what proportion
of this field shares the repetend character of $1/n$? Some elements
will have the same repeating digit or block; others will differ; still
others will terminate altogether. The resulting proportion is a
natural measure of \emph{coherence}: the degree to which the
arithmetic of $n$ organizes its fractional field into a single periodic
pattern.

In this note we compute this measure exactly for a natural family of
denominators, obtain a closed-form alignment formula, and discover that
the golden ratio plays a structural role in the result. (This note
predates the collision invariant series~\cite{nfield} and uses the
earlier repetend/alignment notation rather than the digit-function
conventions of that series.)
The central theorems depend only on the alignment formula
$\alpha = (2m{-}1)/(pm{-}1)$ and the algebra of the golden ratio, and
are therefore independent of the choice of positional base. Decimal
arithmetic provides the motivating instance, but the algebraic content
is base-free.

% ============================================================
\section{Definitions}

We work in base~$10$ throughout. For positive integers $k$ and $n$ with
$1 \le k \le n-1$, the decimal expansion of $k/n$ takes the form
\[
\frac{k}{n} = 0.a_{1}a_{2}\cdots a_{s}\,
\overline{d_{1}d_{2}\cdots d_{\ell}},
\]
where $a_{1}\cdots a_{s}$ is the non-repeating prefix and
$d_{1}\cdots d_{\ell}$ is the repetend of length $\ell \ge 0$. When
$\ell = 0$ the expansion terminates.

\begin{definition}
A positive integer is called \emph{$\{2,5\}$-smooth} (or
\emph{$10$-smooth}) if its only prime factors belong to $\{2,5\}$.
Equivalently, $m$ is $10$-smooth if and only if $1/m$ has a terminating
decimal expansion. We write $\mathcal{S}_{10}$ for the set of all
$10$-smooth positive integers.
\end{definition}

\begin{definition}
The \emph{repetend alignment} of $n$ in base $b$ is the mean,
over $k \in \{1, \ldots, n{-}1\}$, of the position-wise
digit-match proportion between the repetend of $k/n$ and the
repetend of $1/n$:
\[
\alpha(n) = \frac{1}{n-1}
\sum_{k=1}^{n-1} a(k/n,\, 1/n),
\]
where $a(k/n,\, 1/n)$ is defined as follows: if $k/n$
terminates, $a = 1$; if $k/n$ has a repetend of the same length
as $1/n$, $a$ is the fraction of positions at which the two
repetends share the same digit; otherwise $a = 0$. In the family
$n = 3m$ studied here, the repetend of $1/n$ has length one,
so $a$ reduces to exact single-digit matching.
\end{definition}

\begin{remark}
For the family $n = 3m$ with $m \in \mathcal{S}_{10}$, the
repetend of $1/n$ has length one. In this case, the position-wise
definition reduces to exact digit matching: $a = 1$ if the
repeating digit agrees, $a = 0$ otherwise. The general
position-wise definition is used throughout the series to handle
longer repetends at other primes.
\end{remark}

% ============================================================
\section{The Three Classes}

The following lemma establishes the arithmetic that underlies the
alignment count.

\begin{lemma}\label{lem:three-classes}
Let $n = 3m$ with $m \in \mathcal{S}_{10}$ and $m \ge 1$. For each
$k \in \{1,\ldots,n{-}1\}$, exactly one of the following holds:
\begin{enumerate}
\item[\textup{(A)}]
$k \equiv 0 \pmod{3}$. Then $k/n = (k/3)/m$, and since
$m \in \mathcal{S}_{10}$ the expansion terminates.
\item[\textup{(B)}]
$k \equiv 1 \pmod{3}$. Then $k/n$ has repetend of length one,
equal to the repeating digit of $1/n$.
\item[\textup{(C)}]
$k \equiv 2 \pmod{3}$. Then $k/n$ has repetend of length one,
equal to the nines complement of the repeating digit of $1/n$.
\end{enumerate}
\end{lemma}

\begin{proof}
Write $k = 3q + r$ with $r \in \{0,1,2\}$. Then
\[
\frac{k}{3m} = \frac{q}{m} + \frac{r}{3m}.
\]
Since $m \in \mathcal{S}_{10}$, the term $q/m$ terminates.
Adding a terminating decimal does not change the repetend.

\medskip\noindent
\emph{Case $r=0$.}
Then $k/(3m) = q/m$, which terminates.

\medskip\noindent
\emph{Case $r=1$.}
The repetend of $k/(3m)$ equals the repetend of
$1/(3m)$, which is the repeating digit of $1/n$.
(Since $\operatorname{ord}_{3}(10) = 1$, the repetend
has length one.)

\medskip\noindent
\emph{Case $r=2$.}
The repetend of $k/(3m)$ equals the repetend of
$2/(3m)$. Since $1/(3m) + 2/(3m) = 1/m$ terminates,
the repeating digits of $1/(3m)$ and $2/(3m)$ are
nines complements: they sum to~$9$.
\end{proof}

\begin{remark}
The nines complement relationship $d^{*} + (9 - d^{*}) = 9$ between
classes (B) and~(C) reflects the identity $1/3 + 2/3 = 1$ and its
decimal manifestation $0.\overline{3} + 0.\overline{6} = 0.\overline{9}
= 1$. This complement symmetry is a special case of the universal
identity $k/n + (n{-}k)/n = 1$, which pairs every element of
$\mathcal{F}(n)$ with its complement. In base~$10$ the digit pairs
$\{0,9\}$, $\{1,8\}$, $\{2,7\}$, $\{3,6\}$, $\{4,5\}$ each sum
to~$9$, and the alignment classes (B) and~(C) correspond precisely to
the complement pair $\{3,6\}$ associated with the prime~$3$.
\end{remark}

% ============================================================
\section{The Alignment Formula}

\begin{theorem}\label{thm:alignment}
Let $n = 3m$ with $m \in \mathcal{S}_{10}$ and $m \ge 1$. Then
\[
\alpha(n) = \frac{2m-1}{3m-1}.
\]
\end{theorem}

\begin{proof}
By Lemma~\ref{lem:three-classes} the aligned fractions comprise
classes~(A) and~(B). The counts are as follows.

\smallskip

\emph{Class~(A):} The multiples of~$3$ in $\{1,\ldots,3m{-}1\}$ are
$3, 6, \ldots, 3(m{-}1)$, giving $m-1$ elements.

\smallskip

\emph{Class~(B):} The integers $k \equiv 1 \pmod{3}$ in
$\{1,\ldots,3m{-}1\}$ are $1, 4, 7, \ldots, 3m{-}2$, giving $m$
elements.

\smallskip

\emph{Class~(C):} The remaining $m$ elements are not aligned.

\smallskip

The total count of aligned elements is $(m{-}1) + m = 2m - 1$ out of
$3m - 1$, and the formula follows.
\end{proof}

\begin{corollary}\label{cor:limit}
As $m \to \infty$ through $\mathcal{S}_{10}$,
\[
\alpha(3m) \;\longrightarrow\; \frac{2}{3}.
\]
\end{corollary}

The convergence is monotone from below. The first few values are
displayed in Table~\ref{tab:values}.

\begin{table}[h]
\centering
\begin{tabular}{rrrl}
\toprule
$m$ & $n = 3m$ & $\alpha(n)$ & Decimal \\
\midrule
1  &   3 & $1/2$    & $0.500000$ \\
2  &   6 & $3/5$    & $0.600000$ \\
4  &  12 & $7/11$   & $0.636364$ \\
5  &  15 & $9/14$   & $0.642857$ \\
8  &  24 & $15/23$  & $0.652174$ \\
10 &  30 & $19/29$  & $0.655172$ \\
16 &  48 & $31/47$  & $0.659574$ \\
32 &  96 & $63/95$  & $0.663158$ \\
100 & 300 & $199/299$ & $0.665552$ \\
256 & 768 & $511/767$ & $0.666232$ \\
\bottomrule
\end{tabular}
\caption{Repetend alignment $\alpha(3m)$ for selected $10$-smooth
values of $m$.}
\label{tab:values}
\end{table}

% ============================================================
\section{The Golden Threshold}

We now ask: for which $10$-smooth integers $m$ does the alignment exceed
the reciprocal of the golden ratio $\varphi = (1+\sqrt{5})/2$?

\begin{theorem}\label{thm:golden}
Let $m \ge 1$ be a positive integer. Then
\[
\frac{2m-1}{3m-1} \ge \frac{1}{\varphi}
\qquad\Longleftrightarrow\qquad
m \ge \varphi^{2}.
\]
Since $\varphi^{2} = \varphi + 1 = (3+\sqrt{5})/2 \approx 2.618$,
the integer threshold is $m \ge 3$.
\end{theorem}

\begin{proof}
The inequality $(2m{-}1)/(3m{-}1) \ge 1/\varphi$ is equivalent to
\[
\varphi(2m-1) \ge 3m-1.
\]
Expanding and collecting terms:
\[
m(2\varphi - 3) \ge \varphi - 1.
\]
Since $\varphi - 1 = 1/\varphi$ and $\varphi^{2} = \varphi + 1$, we
compute
\[
2\varphi - 3 = 2\varphi - 3 = \frac{2(1+\sqrt{5}) - 6}{2}
= \frac{2\sqrt{5} - 4}{2} = \sqrt{5} - 2.
\]
This is positive (since $\sqrt{5} > 2$), so we may divide:
\[
m \ge \frac{1/\varphi}{\sqrt{5}-2}
= \frac{1}{\varphi(\sqrt{5}-2)}.
\]
Now $\varphi(\sqrt{5}-2) = \varphi\sqrt{5} - 2\varphi$.
Using $\varphi = (1{+}\sqrt{5})/2$:
\[
\varphi\sqrt{5} = \frac{(1+\sqrt{5})\sqrt{5}}{2}
= \frac{\sqrt{5}+5}{2},
\qquad
2\varphi = 1+\sqrt{5},
\]
so
\[
\varphi(\sqrt{5}-2)
= \frac{\sqrt{5}+5}{2} - (1+\sqrt{5})
= \frac{\sqrt{5}+5 - 2 - 2\sqrt{5}}{2}
= \frac{3-\sqrt{5}}{2}.
\]
Therefore
\[
m \ge \frac{2}{3-\sqrt{5}} = \frac{2(3+\sqrt{5})}{(3-\sqrt{5})(3+\sqrt{5})}
= \frac{2(3+\sqrt{5})}{4} = \frac{3+\sqrt{5}}{2} = \varphi^{2}.
\qedhere
\]
\end{proof}

The identity $\varphi^{2} = \varphi + 1$ is the defining property of the
golden ratio~\cite{koshy}. Its appearance here is not assumed but
\emph{derived}: it emerges from the interaction between the golden
threshold $1/\varphi$ and the ternary structure of residue classes
modulo~$3$.

\begin{corollary}\label{cor:classification}
Among denominators of the form $n = 3m$ with $m \in \mathcal{S}_{10}$,
the repetend alignment $\alpha(n)$ exceeds $1/\varphi$ if and only if
$m \ge 4$ \textup{(}the smallest $10$-smooth integer at least
$3$\textup{)}.
\end{corollary}

\begin{proof}
By Theorem~\ref{thm:golden} the integer threshold is $m \ge 3$. The
value $m = 3$ is not $10$-smooth (since $3 \notin \{2,5\}$), so the
smallest admissible value is $m = 4 = 2^{2}$.
\end{proof}

The first admissible denominator is therefore $n = 12$, and
$\alpha(12) = 7/11 \approx 0.6364$, above
$1/\varphi \approx 0.6180$.

\subsection*{The qualifying families}

Two families of integers are proved to satisfy
$\alpha(n) \ge 1/\varphi$:

\begin{enumerate}
\item[\textbf{Tier~1.}]
$n \in \mathcal{S}_{10}$. Every fraction $k/n$ terminates, so
$\alpha(n) = 1$.

\item[\textbf{Tier~2.}]
$n = 3m$ with $m \in \mathcal{S}_{10}$ and $m \ge 4$. Alignment
satisfies $1/\varphi \le \alpha(n) < 1$ by
Theorem~\ref{thm:alignment} and Theorem~\ref{thm:golden}.
\end{enumerate}

Exhaustive computation up to $n \le 909$ finds no integer outside
these two families satisfying $\alpha(n) \ge 1/\varphi$. The
following conjecture asserts that the two families above form the
complete classification.

\begin{conjecture}\label{conj:complete}
For every positive integer $n$, the repetend alignment satisfies
$\alpha(n) \ge 1/\varphi$ if and only if $n \in \mathcal{S}_{10}$
or $n = 3m$ with $m \in \mathcal{S}_{10}$ and $m \ge 4$.
\end{conjecture}

\emph{Note added April 2026: Conjecture~\ref{conj:complete} was
subsequently proved by the author as the three-tier theorem
(``The Three-Tier Theorem,'' research note, 2022). The expected
mechanism, sketched below, is the proof that paper uses.}

The expected mechanism is the nines complement: non-terminating
fractions $k/n$ and $(n{-}k)/n$ have complementary digits at every
position, so their alignment contributions cannot overlap, which
bounds $\alpha(n) \le (t{+}1)/(2t)$ for rough part~$t$. This bound
falls below $1/\varphi$ for $t \ge 7$.

% ============================================================
\section{Why the Golden Ratio, and Why the Prime Three}

The result of Theorem~\ref{thm:golden} invites the question: is there
something special about the threshold $1/\varphi$, or would any
threshold between $3/5$ and $7/11$ produce a similar classification?

A threshold $\tau$ separates the family $\{3m : m \in \mathcal{S}_{10}\}$
into aligned and non-aligned members if and only if it lies in the gap
between $\alpha(6) = 3/5$ and $\alpha(12) = 7/11$. Since
$3/5 < 1/\varphi < 7/11$, the golden ratio does lie in this gap.
In this narrow sense, any $\tau \in (3/5,\, 7/11)$ would work.

What distinguishes $1/\varphi$ is the algebraic structure of the
threshold equation. For a general prime $p$ not dividing the base $b$,
with $b \equiv 1 \pmod{p}$ (so that $\operatorname{ord}_{p}(b) = 1$
and the repetend has length one), the alignment formula generalizes to
\[
\alpha(pm) = \frac{2m-1}{pm-1},
\]
since the classes $k \equiv 0$ and $k \equiv 1 \pmod{p}$ are aligned
while the remaining $p{-}2$ classes are not. The threshold equation
$\alpha(pm) \ge 1/\varphi$ then requires
$m(2\varphi - p) \ge \varphi - 1$.
When $p < 2\varphi \approx 3.236$, the coefficient
$2\varphi - p$ is positive, giving a finite threshold
$m \ge (\varphi - 1)/(2\varphi - p)$.
For $p = 2$ this gives $m \ge 1/2$, trivially satisfied.
For $p = 3$ this gives $m \ge \varphi^2$.

When $p \ge 5$, the coefficient $2\varphi - p$ is negative,
so dividing reverses the inequality:
$m \le (\varphi - 1)/(2\varphi - p) < 0$.
No positive integer satisfies this. Alignment never
reaches $1/\varphi$.

The prime $p = 3$ is therefore the \emph{unique} value for which the
golden threshold produces a non-trivial classification, and the critical
resolution $\varphi^{2} = \varphi + 1$ is the signature of this
uniqueness.

% ============================================================
\section{The Self-Referential Cubic and the Golden Field}

The preceding section establishes that $1/\varphi$ is the unique
golden threshold selecting exactly one non-trivial prime. We now show
that the golden ratio has a deeper algebraic role in the
self-referential structure of the alignment formula.

\begin{definition}
For a threshold $\tau \in (0,1)$ and a prime $p$, the \emph{critical
resolution} is the value $m^{*}(\tau, p)$ at which the alignment
$(2m{-}1)/(pm{-}1)$ equals~$\tau$:
\[
m^{*}(\tau,p) = \frac{1-\tau}{2-p\tau}.
\]
We say the classification is \emph{self-referential} at $(\tau,p)$ if
the critical resolution equals the reciprocal square of the threshold:
\[
m^{*}(\tau,p) = \frac{1}{\tau^{2}}.
\]
\end{definition}

The self-referential condition requires $\tau$ to satisfy the cubic
\begin{equation}\label{eq:cubic}
\tau^{3} - \tau^{2} - p\tau + 2 = 0.
\end{equation}
Since $f(\tau) = \tau^{3} - \tau^{2} - p\tau + 2$ satisfies
$f(0) = 2 > 0$ and $f(1) = 2 - p < 0$ for every prime $p \ge 3$,
the intermediate value theorem guarantees a root in $(0,1)$ for every
such~$p$. The self-referential condition therefore admits a solution
for every prime $p \ge 3$.

What distinguishes $p = 3$ is not the existence of a solution but
its algebraic nature. The following theorem shows that the
self-referential cubic factors over $\mathbb{Q}$ for exactly three
primes, and the two non-trivial cases both produce thresholds in the
golden ratio's number field.

\begin{theorem}\label{thm:cubic-factors}
The self-referential cubic~\eqref{eq:cubic} factors over
$\mathbb{Q}$ if and only if $p \in \{2, 3, 5\}$. The factorizations
are:
\begin{align}
p = 2: \quad & (\tau - 1)(\tau^{2} - 2) = 0, \label{eq:p2} \\
p = 3: \quad & (\tau - 2)(\tau^{2} + \tau - 1) = 0, \label{eq:p3} \\
p = 5: \quad & (\tau + 2)(\tau^{2} - 3\tau + 1) = 0. \label{eq:p5}
\end{align}
For $p \ge 7$, the cubic is irreducible over $\mathbb{Q}$.
\end{theorem}

\begin{proof}
By the rational root theorem, the only possible rational roots of
$\tau^{3} - \tau^{2} - p\tau + 2$ are $\pm 1$ and $\pm 2$.
Evaluating:
\[
f(1) = 2 - p, \quad f(-1) = p, \quad f(2) = 6 - 2p, \quad
f(-2) = 2p - 10.
\]
These vanish at $p = 2, 3, 5$ respectively ($f(-1) = p$ is never
zero for a prime). For $p \ge 7$, no rational root exists and the
cubic is irreducible.
\end{proof}

The case $p = 2$ yields $\tau = 1$ (boundary of the interval) and
$\tau = \sqrt{2}$ (outside $(0,1)$), so it produces no
self-referential threshold. The two non-trivial factorizations are
$p = 3$ and $p = 5$.

\begin{theorem}\label{thm:golden-field}
The self-referential thresholds for $p = 3$ and $p = 5$ are
\[
\tau(3) = \frac{\sqrt{5}-1}{2} = \frac{1}{\varphi},
\qquad
\tau(5) = \frac{3 - \sqrt{5}}{2} = \frac{1}{\varphi^{2}}.
\]
Both lie in the quadratic number field
$\mathbb{Q}(\sqrt{5})$, the field generated by the golden ratio.
They are consecutive powers of $1/\varphi$.

For $p \ge 7$, the self-referential threshold is a cubic irrational
and does not lie in any quadratic number field.
\end{theorem}

\begin{proof}
The quadratic factors in~\eqref{eq:p3} and~\eqref{eq:p5} have
discriminants $1 + 4 = 5$ and $9 - 4 = 5$ respectively, so both
roots lie in $\mathbb{Q}(\sqrt{5})$. For $p = 3$, the root in
$(0,1)$ is $(\sqrt{5}-1)/2 = 1/\varphi$. For $p = 5$, it is
$(3-\sqrt{5})/2$. Since $(3-\sqrt{5})/2 = ((\sqrt{5}-1)/2)^{2}$
(verified by expanding), this equals $1/\varphi^{2}$.

For $p \ge 7$, the cubic is irreducible over $\mathbb{Q}$ by
Theorem~\ref{thm:cubic-factors}, so its roots generate a cubic
extension of $\mathbb{Q}$ and cannot lie in any quadratic field.
\end{proof}

\begin{corollary}\label{cor:fourth-power}
The golden ratio's minimal polynomial $\tau^{2} + \tau - 1$ divides
the self-referential cubic if and only if $p = 3$. The resulting
fourth-power identity is
\[
\tau^{4} = 2 - 3\tau,
\]
where the coefficient~$3$ is the unique prime annihilating the
remainder of the polynomial division.
\end{corollary}

\begin{proof}
Dividing $\tau^{3} - \tau^{2} - p\tau + 2$ by $\tau^{2} + \tau - 1$
yields quotient $\tau - 2$ and remainder $(3 - p)\tau$. The remainder
vanishes if and only if $p = 3$. The fourth-power identity follows
from $\tau^{2} = 1 - \tau$.
\end{proof}

The self-referential structure therefore reveals a hierarchy among
primes:

\begin{enumerate}
\item $p = 3$: threshold $1/\varphi$, governed by the golden ratio's
      own minimal polynomial. The alignment limit $2/3$ exceeds this
      threshold, producing the classification of
      Theorem~\ref{thm:golden}.

\item $p = 5$: threshold $1/\varphi^{2}$, governed by the companion
      quadratic $\tau^{2} - 3\tau + 1$. The alignment limit $2/5$
      exceeds this threshold, but falls below $1/\varphi$, so
      no golden classification arises.

\item $p \ge 7$: thresholds are cubic irrationals outside
      $\mathbb{Q}(\sqrt{5})$. No connection to the golden ratio
      persists.
\end{enumerate}

The golden ratio's number field $\mathbb{Q}(\sqrt{5})$
controls exactly two primes in the self-referential hierarchy.

% ============================================================
\section{Remarks}

\subsection*{Base independence}

The alignment formula $\alpha = (2m{-}1)/(pm{-}1)$ and the
self-referential characterization of Theorem~\ref{thm:golden-field} are
purely algebraic and make no reference to a positional base. The base
enters only in determining which prime~$p$ governs the periodicity of
$1/n$. In any base $b$ satisfying $b \equiv 1 \pmod{3}$
(such as $4$, $7$, $10$, $13$, $16$, or~$19$), the
prime~$3$ plays this role and the full theorem applies with
``$10$-smooth'' replaced by ``$b$-smooth.'' In bases where $3$
divides~$b$ (such as $6$, $12$, or~$30$), a different prime governs
periodicity. Since $p = 3$ is the unique prime selected by the golden
ratio's algebra, the theorem does not apply to such bases. The result
is therefore base-free in its algebraic content and base-dependent only
in whether its hypotheses are met.

\subsection*{Terminating fractions as aligned}

The decision to count terminating fractions as aligned is not
mathematically forced. If terminating fractions are instead excluded, the
alignment among purely periodic fractions is exactly $1/2$ for every $m$,
and the golden ratio plays no role. The classification theorem therefore
reflects a specific but natural convention: that the absence of a
repeating obstruction constitutes compatibility.

\subsection*{Connection to Midy's theorem}

Midy's theorem~\cite{midy,leavitt} and its generalizations study
digit-sum relations within a single repetend, typically for prime
denominators. The present result addresses a complementary question:
not the internal structure of one repetend, but the collective
organization of an entire fractional field. The nines complement
relationship between classes~(B) and~(C) is, however, a manifestation
of the same arithmetic that underlies Midy's observation.

\subsection*{Computational origin}

This result was discovered through computational exploration of the
fractional fields $\mathcal{F}(n)$ for $n \le 909$, using software that
computes structural invariants of decimal expansions~\cite{nfield}. The
empirical observation that a $1/\varphi$ threshold cleanly separates
a family of denominators was observed computationally; the algebraic proof came later.

% ============================================================
\begin{thebibliography}{9}

\bibitem{midy}
E.~Midy,
\emph{De quelques propri\'et\'es des nombres et des fractions
d\'ecimales p\'eriodiques},
Nantes, 1836.

\bibitem{leavitt}
W.~G.~Leavitt,
A theorem on repeating decimals,
\emph{Amer.\ Math.\ Monthly} \textbf{74} (1967), no.~6, 669--673.

\bibitem{hardy}
G.~H.~Hardy and E.~M.~Wright,
\emph{An Introduction to the Theory of Numbers},
6th ed., Oxford University Press, 2008.

\bibitem{koshy}
T.~Koshy,
\emph{Fibonacci and Lucas Numbers with Applications},
2nd ed., Wiley, 2019.

\bibitem{nfield}
A.~S.~Petty,
\emph{nfield: A structural analysis engine for fractional field invariants},
software.
\url{https://github.com/alexspetty/nfield}

\end{thebibliography}

\end{document}
